\documentclass[11pt,a4paper]{article}
\usepackage[margin=27mm]{geometry}
\usepackage{amsmath,amssymb,amsthm,mathtools}
\usepackage{microtype}
\usepackage[hidelinks]{hyperref}
\usepackage{enumitem}
\usepackage{fancyhdr}
\pagestyle{fancy}
\fancyhf{}
\fancyhead[L]{iVAE: almost-everywhere affine identifiability}
\fancyfoot[C]{\thepage}
\setlength{\headheight}{14pt}
\setlength{\parskip}{4pt}
\setlength{\parindent}{0pt}
\newtheorem{theorem}{Theorem}
\newcommand{\R}{\mathbb R}
\newcommand{\Zset}{\mathcal Z}
\newcommand{\Xset}{\mathcal X}
\newcommand{\Uset}{\mathcal U}
\newcommand{\aev}{\text{-a.e.}}
\newcommand{\push}{_{\#}}
\newcommand{\GL}{\operatorname{GL}}
\newcommand{\Leb}{\mathcal L}
\newcommand{\ip}[2]{\langle #1,#2\rangle}
\begin{document}
\begin{center}
{\Large Almost-Everywhere Affine Identifiability of iVAE}\\[5pt]
{\normalsize A measure-theoretic refinement of Theorem 1}
\end{center}

\section*{Notation and Model Setup}
Let $1\le n\le d$, $m,k\ge1$, $q=nk$, and $s\in\{\theta,\theta'\}$.
Let $\Zset\subseteq\R^n$ and $\Xset\subseteq\R^d$ be Borel,
$\Uset\subseteq\R^m$, and
\[
 T_s(z)=\bigl(T_{s,i,j}(z_i)\bigr)_{i,j}\in\R^q,\qquad
 \lambda_s(u)=\bigl(\lambda_{s,i,j}(u)\bigr)_{i,j}\in\R^q.
\]
Put $B_s(z)=\prod_iQ_{s,i}(z_i)$, $Z_s(u)=\prod_iZ_{s,i}(u)$, and
\begin{align}
 p_{s,u}(z)
 &=\frac{B_s(z)}{Z_s(u)}e^{\ip{T_s(z)}{\lambda_s(u)}} \label{eq:prior}\\
 &=\prod_{i=1}^n
 \left[\frac{Q_{s,i}(z_i)}{Z_{s,i}(u)}
 e^{\sum_{j=1}^kT_{s,i,j}(z_i)\lambda_{s,i,j}(u)}\right].
 \nonumber
\end{align}
For injective $f_s:\R^n\to\R^d$ with $f_s(\Zset)=\Xset$, write
$g_s=(f_s|_{\Zset})^{-1}$. Define
\[
 \begin{gathered}
 P_{s,u}(dz)=p_{s,u}(z)\,d\Leb^n|_{\Zset}(z),\qquad t_s=T_s\circ g_s,\\
 \nu_{s,u}=(f_s|_{\Zset})\push P_{s,u},\qquad
 \bar\nu_{s,u}=\iota\push\nu_{s,u},\qquad
 \iota:\Xset\hookrightarrow\R^d.
 \end{gathered}
\]
For a common independent additive noise law $\eta$, let
\begin{equation}
 O_{s,\eta}(u)=\bar\nu_{s,u}*\eta. \label{eq:observed}
\end{equation}
For a fixed measure $\mu$ on $\Xset$, set
\[
 \theta\sim_A^\mu\theta'
 \iff \exists A\in\GL_q(\R),\ c\in\R^q:
 \quad t_\theta=A t_{\theta'}+c\quad\mu\aev
\]
This is an equivalence relation, by the identity, inverse, and composition
of affine maps on common full-$\mu$-measure sets.

\begin{theorem}[Corrected a.e. identifiability]\label{thm:main}
Assume the following.
\begin{enumerate}[label=\textnormal{(\arabic*)},leftmargin=*,itemsep=4pt]
\item \textbf{Injective mixing and measurability.}
\textnormal{(paper assumption (ii); explicit measurable-image interface)}
For both models, $f_s:\R^n\to\R^d$ is injective,
$f_s(\Zset)=\Xset$, and $f_s|_{\Zset}:\Zset\to\Xset$ and
$g_s:\Xset\to\Zset$ are measurable.

\item \textbf{Valid product-form priors.}
\textnormal{(paper equation (7); explicit probability-model validity)}
For both models,
\[
 Q_{s,i}(v)\ge0\quad(\forall i,v),\qquad
 0<Z_{s,i}(u)<\infty\quad(\forall i,u),
\]
and, for every $u\in\Uset$, $p_{s,u}$ is measurable with
$\int_{\Zset}p_{s,u}(z)\,dz=1$.

\item \textbf{Full-frequency noise condition.}
\textnormal{(strengthened assumption (i))}
The measure $\eta$ is a probability measure on $\R^d$ and
\begin{equation}
 \Leb^d\{\omega\in\R^d:\varphi_\eta(\omega)=0\}=0,\qquad
 \varphi_\eta(\omega)=\int e^{\mathrm i\ip{\omega}{y}}\,d\eta(y).
 \label{eq:noise}
\end{equation}

\item \textbf{Auxiliary variation.}
\textnormal{(paper assumption (iv))}
There exist pairwise distinct $u_0,\ldots,u_q\in\Uset$ such that
\begin{equation}
 L=[\lambda_\theta(u_1)-\lambda_\theta(u_0)\ \cdots\
      \lambda_\theta(u_q)-\lambda_\theta(u_0)],\qquad \det L\ne0.
 \label{eq:L}
\end{equation}

\item \textbf{Essential affine nondegeneracy of the true pivot prior.}
\textnormal{(additional assumption for Step III; related to assumption (iii))}
For every $v\in\R^q$, $b\in\R$,
\begin{equation}
 [\,\ip{v}{T_\theta(z)}=b\quad P_{\theta,u_0}\aev\,]
 \ \Longrightarrow\ v=0. \label{eq:nd}
\end{equation}
This condition is assumed, not deduced from the paper's assumption (iii).
\end{enumerate}
If
\begin{equation}
 O_{\theta,\eta}(u_r)=O_{\theta',\eta}(u_r)
 \qquad(r=0,\ldots,q), \label{eq:hobs}
\end{equation}
then $\sigma:=\nu_{\theta,u_0}=\nu_{\theta',u_0}$ and
\begin{equation}
 \exists A\in\GL_q(\R),\ c\in\R^q:
 \quad t_\theta(x)=A t_{\theta'}(x)+c\quad\sigma\aev \label{eq:conclusion}
\end{equation}
Consequently,
\begin{equation}
 \mu\ll\sigma\ \Longrightarrow\ \theta\sim_A^\mu\theta'.
 \label{eq:equivalence}
\end{equation}
The same $A,c$ apply to every such $\mu$, including all $\nu_{s,u}$.
\end{theorem}

\begin{proof}\mbox{}\par
\noindent\textbf{Step I.}
For $r=0,\ldots,q$, the convolution identity and \eqref{eq:hobs} give
\[
 \varphi_{\bar\nu_{\theta,u_r}}\varphi_\eta
 =\varphi_{\bar\nu_{\theta',u_r}}\varphi_\eta.
\]
By \eqref{eq:noise}, continuity of characteristic functions, and Fourier
uniqueness for finite measures,
\[
 \begin{aligned}
 &\varphi_{\bar\nu_{\theta,u_r}}
   =\varphi_{\bar\nu_{\theta',u_r}}\quad\Leb^d\aev\\
 &\Longrightarrow\
   \varphi_{\bar\nu_{\theta,u_r}}
   =\varphi_{\bar\nu_{\theta',u_r}}\quad\text{everywhere}\\
 &\Longrightarrow\ \bar\nu_{\theta,u_r}=\bar\nu_{\theta',u_r}.
 \end{aligned}
\]
Here continuity upgrades a.e. equality since every nonempty open subset
of $\R^d$ has positive Lebesgue measure.
The inclusion $\iota$ is a measurable embedding; hence
\begin{equation}
 \nu_{\theta,u_r}=\nu_{\theta',u_r}\qquad(r=0,\ldots,q).
 \label{eq:hlaw}
\end{equation}
In particular, $\sigma=\nu_{\theta,u_0}=\nu_{\theta',u_0}$ is a probability measure.

\medskip
\noindent\textbf{Step II.}
Fix $s\in\{\theta,\theta'\}$ and $u\in\Uset$. By \eqref{eq:prior},
\begin{equation}
 p_{s,u}(z)=0\iff B_s(z)=0\iff p_{s,u_0}(z)=0. \label{eq:zeros}
\end{equation}
Define the nonnegative measurable ratio
\begin{equation}
 r_{s,u}(z)=
 \begin{cases}
 p_{s,u}(z)/p_{s,u_0}(z),&p_{s,u_0}(z)>0,\\
 0,&p_{s,u_0}(z)=0.
 \end{cases} \label{eq:ratio}
\end{equation}
Using \eqref{eq:zeros},
\begin{equation}
 p_{s,u_0}r_{s,u}=p_{s,u}\quad\text{on }\Zset, \label{eq:densityproduct}
\end{equation}
and therefore
\begin{equation}
 P_{s,u}=r_{s,u}\,P_{s,u_0}. \label{eq:latentweight}
\end{equation}
Let $R_{s,u}=r_{s,u}\circ g_s$. For every measurable $E\subseteq\Xset$,
\[
 \nu_{s,u}(E)
 =\int_{(f_s|_{\Zset})^{-1}(E)}r_{s,u}\,dP_{s,u_0}
 =\int_E R_{s,u}\,d\nu_{s,u_0}.
\]
Thus
\begin{equation}
 \nu_{s,u}=R_{s,u}\,\nu_{s,u_0},\qquad
 \nu_{s,u}\ll\nu_{s,u_0}. \label{eq:imageweight}
\end{equation}
Since $P_{s,u_0}\{p_{s,u_0}=0\}=0$, cancellation in \eqref{eq:ratio} gives
\[
 r_{s,u}(z)=\frac{Z_s(u_0)}{Z_s(u)}
 e^{\ip{T_s(z)}{\lambda_s(u)-\lambda_s(u_0)}}>0
 \quad P_{s,u_0}\aev
\]
Set $a_s(u)=\log Z_s(u)$ and
$\Delta\lambda_s(u)=\lambda_s(u)-\lambda_s(u_0)$. Taking logarithms and
pushing forward,
\begin{equation}
 \log R_{s,u}(x)=\ip{t_s(x)}{\Delta\lambda_s(u)}
 -[a_s(u)-a_s(u_0)]\quad\nu_{s,u_0}\aev \label{eq:logratio}
\end{equation}
From \eqref{eq:hlaw} and \eqref{eq:imageweight},
\[
 R_{\theta,u_r}\,\sigma=R_{\theta',u_r}\,\sigma.
\]
Uniqueness of densities with respect to the same finite measure yields
\begin{equation}
 R_{\theta,u_r}=R_{\theta',u_r}\quad\sigma\aev \label{eq:ratioeq}
\end{equation}
Combining \eqref{eq:logratio}--\eqref{eq:ratioeq},
\begin{align}
 \ip{t_\theta(x)}{\Delta\lambda_\theta(u_r)}
 -[a_\theta(u_r)-a_\theta(u_0)]
 &=\ip{t_{\theta'}(x)}{\Delta\lambda_{\theta'}(u_r)}\nonumber\\
 &\quad-[a_{\theta'}(u_r)-a_{\theta'}(u_0)]
 \quad\sigma\aev \label{eq:pivotrow}
\end{align}
Intersecting finitely many full-$\sigma$-measure sets makes these identities
simultaneous for $r=1,\ldots,q$. With
\[
 \begin{aligned}
 L'&=[\Delta\lambda_{\theta'}(u_1)\ \cdots\
          \Delta\lambda_{\theta'}(u_q)],\\
 b_r&=[a_\theta(u_r)-a_\theta(u_0)]
       -[a_{\theta'}(u_r)-a_{\theta'}(u_0)],
 \end{aligned}
\]
equation \eqref{eq:pivotrow} becomes
\begin{equation}
 L^\top t_\theta=(L')^\top t_{\theta'}+b\quad\sigma\aev \label{eq:matrix}
\end{equation}
Since $\det L\ne0$, set
\begin{equation}
 A=(L^\top)^{-1}(L')^\top,\qquad c=(L^\top)^{-1}b. \label{eq:witness}
\end{equation}
Then
\begin{equation}
 t_\theta=A t_{\theta'}+c\quad\sigma\aev \label{eq:stepii}
\end{equation}

\medskip
\noindent\textbf{Step III.}
Using $\sigma=(f_\theta|_{\Zset})\push P_{\theta,u_0}$ and
$t_\theta\circ f_\theta|_{\Zset}=T_\theta|_{\Zset}$, assumption \eqref{eq:nd} gives
\[
 \begin{aligned}
 [\,\ip{v}{t_\theta}=b\quad\sigma\aev\,]
 &\Longrightarrow
 [\,\ip{v}{T_\theta}=b\quad P_{\theta,u_0}\aev\,]\\
 &\Longrightarrow v=0.
 \end{aligned}
\]
If $A$ were singular, choose $0\ne v\in\ker A^\top$. By \eqref{eq:stepii},
\[
 \ip{v}{t_\theta}
 =\ip{A^\top v}{t_{\theta'}}+\ip{v}{c}
 =\ip{v}{c}\quad\sigma\aev,
\]
forcing $v=0$, a contradiction. Thus $A\in\GL_q(\R)$.

For $\mu\ll\sigma$, the exceptional set in \eqref{eq:stepii} is also
$\mu$-null, proving \eqref{eq:equivalence} with the same $A,c$.
Finally, \eqref{eq:imageweight} implies
$\nu_{s,u}\ll\nu_{s,u_0}=\sigma$ for every $s,u$.
\end{proof}

\noindent\emph{Scope.}
Only injectivity on $\Zset$ is used; no ambient surjectivity or
$n=d$ assumption is required. The conclusion is a.e., not pointwise.
The full-frequency condition and \eqref{eq:nd} are explicit corrections
to the original theorem of \cite{ivae}.

\begin{thebibliography}{1}
\bibitem{ivae}
I.~Khemakhem, D.~P.~Kingma, R.~Monti, and A.~Hyv\"arinen.
\emph{Variational Autoencoders and Nonlinear ICA: A Unifying Framework}.
arXiv:1907.04809v4, 2020. Theorem 1 and Appendix B.2.
\url{https://arxiv.org/abs/1907.04809v4}.
\end{thebibliography}
\end{document}
